% ---
% Reemplaza líneas 2881–2890 de Docs/Technical_Report.tex (y crea sección propia M6).
% Decisiones que lo respaldan: D-016, D-017, D-027, D-028, D-030, D-032, D-033, D-042,
% H-013, H-032, H-033, H-037, H-038, H-046, H-055.
% ---

\subsection{Módulo 6: Análisis de separabilidad}
\label{sec:modulo6}

El análisis de separabilidad opera sobre dos objetos distintos, y conviene
distinguirlos al interpretar los resultados. Todas las métricas se calculan sobre
la misma submuestra: 200 lesiones del conjunto de entrenamiento por subconjunto,
estratificadas por clase y elegidas con la semilla 42; el conjunto de prueba no interviene.
El tamaño se fijó antes de ver los datos por el coste del protocolo físico: en un equipo cuántico,
cada entrada del kernel es un circuito \textit{compute-uncompute} distinto, de modo que una matriz de
200 lesiones requiere 19\,900 circuitos, frente a cerca de 0.87 y 1.19 millones con todas las masas y
todas las calcificaciones de entrenamiento. En simulación, el cálculo a partir de los vectores de estado
hace viable el conjunto completo; por ello, el experimento X9 (sección~\ref{sec:experimentos_complementarios})
repite las métricas sobre otras 19 submuestras y sobre todo el entrenamiento como análisis de robustez.
Los kernels cuánticos se evalúan de forma exacta a partir de los vectores de estado, y la
significancia de cada métrica que usa las etiquetas se estima con 1\,000 permutaciones de las etiquetas que
conservan el número de lesiones de cada clase.

\textbf{Métricas sobre kernels.} El kernel cuántico de fidelidad
$K_Q(\mathbf{x},\mathbf{x}') = |\langle\phi(\mathbf{x})|\phi(\mathbf{x}')\rangle|^2$ se
evalúa de forma exacta sobre el \textit{feature map} exclusivamente, sin \textit{ansatz}.
Se calculan tres cantidades:
\begin{itemize}
  \item El \textit{kernel-target alignment} centrado \cite{Cortes2012}, con la versión sin
        centrar \cite{cristianini2002kernel} como complemento. Dado que para un kernel cercano a la
        identidad ambas valen aproximadamente $1/\sqrt{N}$ con independencia de las
        etiquetas, se compara el exceso sobre la media de una prueba de permutación.
  \item La dimensión efectiva $d$ del espacio que generan las lesiones \cite{huang2021power}
        (ecuación~\eqref{eq:dim_efectiva}, sección~\ref{sec:metricas}), que cuantifica en cuántas
        direcciones independientes se reparten las lesiones dentro del espacio del kernel: va de
        1, si todas se parecen, a $N$, si cada lesión es ortogonal a las demás, calculada a partir
        de los valores propios ordenados de la matriz del kernel normalizada.
  \item La \textit{geometric difference} regularizada $g_{\mathrm{gen}}$ \cite{huang2021power}
        (ecuaciones~\eqref{eq:g_gen} y~\eqref{eq:g_tra}, sección~\ref{sec:metricas}), sujeta a que
        el error de entrenamiento $g_{\mathrm{tra}}$ sea menor que 0.045.
\end{itemize}

El kernel clásico de referencia $K_C$ es el RBF con el ancho de banda $\gamma$ fijado por la
heurística de la mediana, el mismo de C3. Siguiendo el protocolo de Huang et al. \cite{huang2021power},
se reporta el mínimo sobre una batería de kernels clásicos (un kernel lineal sobre los datos
centrados y nueve kernels gaussianos con $\gamma = m/(n\operatorname{Var}[\mathbf{x}])$, para
$m\in\{0.25,0.5,1,2,4,8,16,32,64\}$) y valores de regularización
$\lambda\in\{10^{-5},10^{-4},10^{-3},10^{-2},0.025,0.05,0.1\}$, con $\operatorname{Tr}(K) = N$. Como
sensibilidad se añaden tres gaussianos más estrechos, $m\in\{128,256,512\}$;
para cada kernel se toma el mayor $\lambda$ admisible que satisface la cota sobre $g_{\mathrm{tra}}$.
Sin regularización, la cantidad mide sobre todo el condicionamiento de $K_C$ y no la codificación.

Como análisis de sensibilidad, los kernels cuánticos se recalculan con los ángulos escalados,
$\mathbf{x}\to c\,\mathbf{x}$, para $c \in \{1, 0.5, 0.2, 0.1, 0.05, 0.02\}$. Así se mide cómo cambia la
geometría cuando las rotaciones son menores.

\textbf{Métricas sobre los \textit{embeddings}.} El índice de Davies-Bouldin, la razón de
Fisher multivariante (ecuación~\eqref{eq:fisher_multivariate}) y el alineamiento del kernel lineal se
calculan sobre los vectores de las cinco condiciones, que son los que recibe el
clasificador; la visualización \textit{t-SNE} (perplejidad 30, inicialización por PCA,
semilla 42) se usa solo como apoyo visual. Las tres métricas son invariantes ante la
rotación de C2, que actúa así como control nulo.

Ambas familias de métricas no son equivalentes: el kernel de fidelidad compara estados
completos, incluida su información de fase, mientras que los valores esperados en la
base $Z$ solo ven la parte de esa información que el \textit{ansatz} convierte en probabilidades.
Por eso las dos familias pueden ordenar las condiciones de forma distinta.
